\chapter*{Anexo A: Instalación, entorno y comandos principales}
\addcontentsline{toc}{chapter}{Anexo A: Instalación, entorno y comandos principales}
\label{cap:anexo_instalacion}

\setcounter{section}{0}
\renewcommand{\thesection}{A.\arabic{section}}
\setcounter{code}{0}
\renewcommand{\thecode}{A.\arabic{code}}

Este anexo recoge, de forma reproducible, los pasos necesarios para preparar
el entorno de desarrollo, regenerar los artefactos derivados del paquete
oficial CVN y ejecutar la herramienta local de gestión descrita en el
Capítulo~\ref{cap:formato_herramienta}. No introduce ninguna decisión de
diseño nueva: es una referencia operativa que complementa, con la sintaxis
exacta de los comandos, la Tabla~\ref{tab:formato_comandos}.

\section{Prerrequisitos y configuración del entorno}
\label{sec:anexo_instalacion_entorno}

El proyecto requiere \texttt{uv} como gestor de entorno y dependencias, y
Python 3.14. El Código~\ref{cod:anexo_entorno} resume la secuencia completa
de preparación: fijar la versión de Python, sincronizar los grupos de
dependencias necesarios e instalar el proyecto en modo editable, requisito
para que el código bajo \texttt{src/} sea importable durante la ejecución de
pruebas y de la herramienta local.

\begin{code}[htbp]
\begin{lstlisting}[language=bash, basicstyle=\footnotesize\ttfamily, breaklines=true]
# Fijar la version de Python del proyecto (opcional si ya esta activa)
uv python pin 3.14

# Sincronizar dependencias de generacion estructural y de pruebas
uv sync --group codegen --group testing

# Instalar el proyecto en modo editable
uv pip install -e .

# Verificar que la herramienta local esta disponible
uv run open-cvn --help
\end{lstlisting}
\caption{Preparación del entorno de desarrollo.}
\label{cod:anexo_entorno}
\end{code}

El grupo de dependencias \texttt{codegen} incluye \texttt{xsdata} y
\texttt{xsdata-pydantic}, empleados por la generación estructural descrita en
la Sección~\ref{sec:pipeline_generacion_estructural}; el grupo
\texttt{testing} incluye \texttt{pytest} y \texttt{pytest-xdist}, empleados
por la batería de pruebas evaluada en el Capítulo~\ref{cap:evaluacion}.

\section{Regeneración de artefactos y verificación}
\label{sec:anexo_instalacion_regeneracion}

Los artefactos generados bajo \texttt{src/generated/} no se editan a mano:
se regeneran a partir de los esquemas XSD oficiales cuando estos cambian. El
Código~\ref{cod:anexo_regeneracion} recoge el comando que regenera los tres
objetivos estructurales (paquete CVN, manual de especificación y modelo en
árbol) y la verificación completa mediante la batería de pruebas
automatizadas, la misma que documenta el Capítulo~\ref{cap:evaluacion}.

\begin{code}[htbp]
\begin{lstlisting}[language=bash, basicstyle=\footnotesize\ttfamily, breaklines=true]
# Regenerar los tres objetivos estructurales
uv run python -m cvn_codegen.xsdata_runner all

# Ejecutar la bateria completa de pruebas en paralelo
uv run pytest -n auto tests
\end{lstlisting}
\caption{Regeneración de artefactos estructurales y verificación por pruebas.}
\label{cod:anexo_regeneracion}
\end{code}

\section{Recorrido completo con la herramienta local}
\label{sec:anexo_instalacion_recorrido}

El Código~\ref{cod:anexo_recorrido} muestra un recorrido de extremo a
extremo con la herramienta local \texttt{open-cvn}: inicialización de un
almacén local, importación de un documento Open CVN JSON como currículo
maestro, creación de una versión derivada, exploración y exclusión de una
entrada, y exportación a Open CVN JSON y a LaTeX. Cada subcomando se
corresponde con una fila de la Tabla~\ref{tab:formato_comandos}.

\begin{code}[htbp]
\begin{lstlisting}[language=bash, basicstyle=\footnotesize\ttfamily, breaklines=true]
STORE=/tmp/open-cvn-demo.sqlite

# 1. Inicializar el almacen local
uv run open-cvn store init --path $STORE

# 2. Importar un documento Open CVN JSON como maestro
uv run open-cvn json import examples/open_cvn/research_entry.json \
  --store $STORE --as-master

# 3. Crear una version derivada
uv run open-cvn versions derive public --store $STORE

# 4. Explorar las secciones de la version derivada
uv run open-cvn versions sections public --store $STORE

# 5. Excluir una entrada concreta de la version derivada
uv run open-cvn versions exclude public /curriculum/research/0 \
  --store $STORE

# 6. Exportar la version derivada a Open CVN JSON
uv run open-cvn json export /tmp/public.json \
  --store $STORE --version public

# 7. Exportar la version derivada a LaTeX
uv run open-cvn latex export /tmp/public.tex \
  --store $STORE --version public
\end{lstlisting}
\caption{Recorrido completo de importación, versionado y exportación con la herramienta local.}
\label{cod:anexo_recorrido}
\end{code}

La generación de PDF a partir de la exportación LaTeX, mediante
\texttt{open-cvn pdf generate}, es opcional y depende de que exista un motor
LaTeX disponible en el sistema (\texttt{latexmk} o \texttt{pdflatex}); en su
ausencia, la herramienta lo reporta de forma explícita en lugar de fallar de
forma silenciosa. La importación de currículos desde PDF mediante modelos de
lenguaje de gran tamaño exige, además, activar explícitamente el indicador
\texttt{--allow-external-llm}, precisamente porque un PDF puede contener
datos personales; este mecanismo se describe con más detalle en la
Sección~\ref{sec:formato_importacion} y en la
Sección~\ref{sec:evaluacion_discusion}.
